% Part: first-order-logic
% Chapter: beyond
% Section: intuitionistic-logic

\documentclass[../../../include/open-logic-section]{subfiles}

\begin{document}

\olfileid[pt]{fol}{byd}{il}

\olsection{Lógica Intuicionista}

Em contraste com a lógica de segunda ordem e de ordem superior, a lógica
intuicionista de primeira ordem representa uma restrição da versão
clássica, destinada a modelar um tipo de raciocínio mais
``construtivo''. Os exemplos a seguir podem ilustrar algumas das
motivações subjacentes.

Suponha que alguém se aproximasse de você um dia e anunciasse que havia
determinado um número natural~$x$, com a propriedade de que se $x$ é
primo, a hipótese de Riemann é verdadeira, e se $x$ é composto, a
hipótese de Riemann é falsa. Ótimas notícias!{} Se a hipótese de Riemann
é verdadeira ou não é uma das grandes questões em aberto da
matemática, e aqui parecem ter reduzido o problema a um cálculo, isto
é, à determinação de se um número específico é primo ou não.

Qual é o valor mágico de $x$? Eles o descrevem assim: $x$ é o número
natural que é igual a $7$ se a hipótese de Riemann é verdadeira, e $9$
caso contrário.

Com raiva, você exige seu dinheiro de volta. De um ponto de vista
clássico, a descrição acima de fato determina um valor único de $x$;
mas o que você realmente quer é um valor de $x$ que seja dado
\emph{explicitamente}.

Para tomar outro exemplo, talvez menos artificial, considere a
seguinte questão. Sabemos que é possível elevar um número irracional a
uma potência racional e obter um resultado racional. Por exemplo,
$\sqrt{2}^2 = 2$. O que é menos claro é se é possível elevar um número
irracional a uma potência \emph{irracional} e obter um resultado
racional. O seguinte teorema responde afirmativamente:

\begin{thm}
Existem números irracionais $a$ e $b$ tais que $a^b$ é racional.
\end{thm}

\begin{proof}
Considere $\sqrt{2}^{\sqrt{2}}$. Se isto é racional, terminamos:
podemos tomar $a = b = \sqrt{2}$. Caso contrário, é irracional. Então
temos
\[
(\sqrt{2}^{\sqrt{2}})^{\sqrt{2}} = \sqrt{2}^{\sqrt{2} \cdot
  \sqrt{2}} = \sqrt{2}^2 = 2,
\]
que é certamente racional. Assim, neste caso, tome $a$ como
$\sqrt{2}^{\sqrt{2}}$, e $b$ como~$\sqrt 2$.
\end{proof}

Isso constitui uma prova válida? A maioria dos matemáticos sente que
sim. Mas, novamente, há algo um pouco insatisfatório aqui: provamos a
existência de um par de números reais com uma certa propriedade, sem
poder dizer \emph{qual} par de números é.  É possível provar o mesmo
resultado, mas de tal forma que o par $a$, $b$ \emph{é} dado na prova:
tome $a = \sqrt{3}$ e $b = \log_3 4$. Então
\[
a^b = \sqrt{3}^{\log_3 4} = 3^{1/2 \cdot \log_3 4} = (3^{\log_3
  4})^{1/2} = 4^{1/2}= 2,
\]
pois $3^{\log_3 x} = x$.

A lógica intuicionista foi concebida para modelar um tipo de
raciocínio no qual movimentos como o da primeira prova são
proibidos. Provar a existência de um $x$ que satisfaz~$!A(x)$ significa
que se deve fornecer um~$x$ específico e uma prova de que ele satisfaz
$!A$, como na segunda prova. Provar que $!A$ ou $!B$ vale requer que
se possa provar um ou outro.

Formalmente falando, a lógica intuicionista de primeira ordem é o que se
obtém se se restringe um sistema de !!{derivation} para a lógica de
primeira ordem de certa maneira. Similarmente, há versões intuicionistas
da lógica de segunda ou de ordem superior. Do ponto de vista matemático,
estes são apenas sistemas dedutivos formais, mas, como já observado, são
destinados a modelar um tipo de raciocínio matemático. Pode-se tomar
isso como o tipo de raciocínio justificado por uma certa visão
filosófica da matemática (como o intuicionismo de Brouwer); pode-se
tomá-lo como um tipo de raciocínio matemático mais ``concreto'' e
satisfatório (na linha do construtivismo de Bishop); e pode-se discutir
se a descrição formal captura ou não a motivação informal. Mas quaisquer
que sejam as posições filosóficas que possamos ter, podemos estudar a
lógica intuicionista como uma lógica formalmente apresentada; e por
quaisquer razões, muitos lógicos matemáticos acham interessante fazê-lo.

Há uma interpretação construtiva informal dos conectivos intuicionistas,
usualmente conhecida como interpretação BHK (em homenagem a Brouwer,
Heyting e Kolmogorov). Ela funciona assim: uma prova de $!A
\land !B$ consiste em uma prova de $!A$ emparelhada com uma prova de
$B$; uma prova de $!A \lor !B$ consiste em uma prova de $!A$, ou uma
prova de $!B$, onde temos informação explícita sobre qual é o caso; uma
prova de $!A \lif !B$ consiste em um procedimento que transforma uma
prova de $!A$ em uma prova de~$!B$; uma prova de $\lforall[x][!A(x)]$
consiste em um procedimento que retorna uma prova de $!A(x)$ para
qualquer valor de~$x$; e uma prova de $\lexists[x][!A(x)]$ consiste em
um valor de~$x$, juntamente com uma prova de que esse valor
satisfaz~$!A$. Pode-se descrever a interpretação em termos
computacionais conhecidos como o ``isomorfismo de Curry--Howard'' ou o
``paradigma !!{formula}s-como-tipos'': pense em !!a{formula} como
especificando um certo tipo de dado, e provas como objetos computacionais
desses tipos de dados que nos permitem ver que a !!{formula} correspondente
é verdadeira.

A lógica intuicionista é frequentemente pensada como a lógica clássica
``menos'' a lei do terceiro excluído. O seguinte teorema torna isso mais
preciso.
\begin{thm}
Intuicionisticamente, os seguintes esquemas de axiomas são equivalentes:
\begin{enumerate}
\item $(\lnot !A \lif \lfalse) \lif !A$.
\item $!A \lor \lnot !A$
\item $\lnot \lnot !A \lif !A$
\end{enumerate}
\end{thm}
Obter instâncias de um esquema a partir de qualquer um dos outros é um bom
exercício em lógica intuicionista.

Os primeiros sistemas dedutivos para a lógica proposicional
intuicionista, apresentados como formalizações do intuicionismo de
Brouwer, são devidos, independentemente, a Kolmogorov, Glivenko e
Heyting. A primeira formalização da lógica intuicionista de primeira
ordem (e partes da matemática intuicionista) é devida a Heyting. Embora
vários esquemas válidos clássicamente não sejam válidos
intuicionisticamente, muitos o são.

A \emph{tradução de dupla negação} descreve uma relação importante
entre a lógica clássica e a intuicionista. Ela é definida
indutivamente da seguinte forma (pense em $!A^N$ como a tradução
``intuicionista'' da !!{formula} clássica~$!A$):
\begin{align*}
!A^N & \ident \lnot\lnot !A \quad \text{para !!{formula}s atômicas $!A$} \\
(!A \land !B)^N & \ident (!A^N \land !B^N) \\
(!A \lor !B)^N & \ident  \lnot\lnot (!A^N \lor !B^N) \\
(!A \lif !B)^N & \ident (!A^N \lif !B^N) \\
(\lforall[x][!A])^N & \ident \lforall[x][!A^N] \\
(\lexists[x][!A])^N & \ident \lnot\lnot\lexists[x][!A^N]
\end{align*}
Kolmogorov e Glivenko tiveram versões dessa tradução para a lógica
proposicional; para a lógica de predicados, é devida a G\"odel e
Gentzen, independentemente. Temos

\begin{thm}
\begin{enumerate}
\item $!A \liff !A^N$ é demonstrável classicamente
\item Se $!A$ é demonstrável classicamente, então $!A^N$ é demonstrável
  intuicionisticamente.
\end{enumerate}
\end{thm}

Podemos agora imaginar o seguinte diálogo. Matemático clássico:
``Provei $!A$!'' Matemático intuicionista: ``Sua prova não é
válida. O que você realmente provou é $!A^N$.'' Matemático clássico:
``Tudo bem para mim!'' Quanto ao matemático clássico, o intuicionista
está apenas sendo minucioso, pois os dois são equivalentes. Mas o
intuicionista insiste que há uma diferença.

Note que a tradução acima diz respeito apenas à lógica pura; ela não
aborda a questão de quais são os axiomas \emph{não lógicos} adequados
para a matemática clássica e intuicionista, ou qual é a relação entre
eles. Mas a seguinte extensão ligeira do teorema acima fornece alguma
informação útil:

\begin{thm}
Se $\Gamma$ prova $!A$ classicamente, então $\Gamma^N$ prova $!A^N$
intuicionisticamente.
\end{thm}

Em outras palavras, se $!A$ é demonstrável a partir de algumas hipóteses
classicamente, então $!A^N$ é demonstrável a partir de suas traduções
de dupla negação.

Para mostrar que uma sentença ou !!{formula} proposicional é
intuicionisticamente válida, basta fornecer uma prova. Mas como se pode
mostrar que ela não é válida? Para esse propósito, precisamos de uma
semântica que seja correta e, de preferência, completa. Uma semântica
devida a Kripke se encaixa bem.

Podemos jogar o mesmo jogo que fizemos para a lógica clássica: definir a
semântica e provar correção e completude. Vale a pena, no entanto, notar
a seguinte distinção. No caso da lógica clássica, a semântica era a
``óbvia'', num sentido implícito no significado dos conectivos. Embora
se possa fornecer alguma motivação intuitiva para a semântica de Kripke,
esta não oferece o mesmo sentimento de inevitabilidade. Além disso, a
noção de !!{structure} clássica é uma noção matemática natural, de modo
que podemos ou tomar a noção de !!{structure} como uma ferramenta para
estudar a lógica clássica de primeira ordem, ou tomar a lógica clássica
de primeira ordem como uma ferramenta para estudar !!{structure}s
matemáticas. Em contraste, as !!{structure}s de Kripke só podem ser
vistas como um construto lógico; elas não parecem ter interesse
matemático independente.

Uma !!{structure} de Kripke~$\mModel M = \tuple{W, R, V}$ para uma
linguagem proposicional consiste em um conjunto~$W$, uma ordem parcial~$R$
em~$W$ com um menor !!{element}, e uma atribuição ``monótona'' de
variáveis proposicionais aos !!{element}s de~$W$. A intuição é que os
!!{element}s de $W$ representam ``mundos'', ou ``estados de conhecimento'';
um elemento $v \geq u$ representa um ``possível estado futuro'' de~$u$; e
as variáveis proposicionais atribuídas a~$u$ são as proposições que se
sabe serem verdadeiras no estado~$u$. A relação de forçamento
$\mSat{M}{!A}[w]$ estende essa relação a !!{formula}s arbitrárias na
linguagem; leia $\mSat{M}{!A}[w]$ como ``$!A$ é verdadeira no
estado~$w$.'' A relação é definida indutivamente, da seguinte forma:
\begin{enumerate}
\item $\mSat{M}{\Obj p_i}[w]$ se, e somente se, $\Obj p_i$ é uma das
  variáveis proposicionais atribuídas a~$w$.
\item $\mSat/{M}{\lfalse}[w]$.
\item $\mSat{M}{(!A \land !B)}[w]$ se, e somente se, $\mSat{M}{!A}[w]$ e $\mSat{M}{!B}[w]$.
\item $\mSat{M}{(!A \lor !B)}[w]$ se, e somente se, $\mSat{M}{!A}[w]$ ou $\mSat{M}{!B}[w]$.
\item $\mSat{M}{(!A \lif !B)}[w]$ se, e somente se, sempre que $w' \geq w$ e
  $\mSat{M}{!A}[w']$, então $\mSat{M}{!B}[w']$.
\end{enumerate}
É um bom exercício tentar mostrar que $\lnot (p \land q) \lif
(\lnot p \lor \lnot q)$ não é intuicionisticamente válida, preparando uma
!!{structure} de Kripke que forneça um contraexemplo.

\end{document}
